\section{混杂、碰撞点与后门调整}
\label{sec:13-Machine-Learning/12-Causal-Inference/03}

一位分析师在估计某药物对住院时长的影响。第一版模型只放了年龄和基础病，效应是 \(-1.2\) 天。审稿意见说协变量太少，于是他把数据表里剩下的列全加了进去：入院科室、是否完成随访、用药后第七天的化验指标、有没有参加满意度回访。效应变成 \(+0.4\) 天。

两版都能跑通，两版的拟合优度后者还更高。他没法判断哪个数字更可信，因为``多控制一些变量''这条直觉在这里给不出方向——它默认所有协变量的角色相同，而这恰恰不成立。控制共同原因关掉一条虚假通路，控制碰撞点却打开一条本来关着的，控制中介则删掉了想要计入的那部分效应。三种变量都可能与结果高度相关，回归看不出它们的区别。

\subsection{后门路径是相关性里那部分不属于因果的东西}

从 \(X\) 出发到 \(Y\) 的路径，如果第一步的箭头是指向 \(X\) 的，称为后门路径。最简单的一条是

\[
X\longleftarrow C\longrightarrow Y,
\]

\(C\) 同时影响处理与结果。即便 \(X\) 对 \(Y\) 毫无作用，观察数据里两者也会相关——相关性沿着这条通路漏进来了。识别的目标就是把所有这类通路堵上，只留下 \(X\to\cdots\to Y\) 的有向路径。

后门准则把这件事写成判据：若一组变量 \(C\) 阻断了 \(X\) 与 \(Y\) 之间的所有后门路径，且 \(C\) 不含 \(X\) 的后代，则

\[
p\bigl(y\given\operatorname{do}(X=x)\bigr)=\sum_c p(y\given x,c)\,p(c).
\]

骨架很短。按截断分解，干预删掉的是 \(p(x\given \operatorname{pa}(X))\) 这个因子；后门条件保证 \(C\) 把 \(X\) 与它的非后代祖先隔开，于是求和后剩下的表达式只含观察分布。条件对账有三条：\(C\) 必须实际阻断每一条后门路径（漏一条就有残余混杂）；\(C\) 不能含 \(X\) 的后代（否则删掉的因子之外还有被污染的项）；每个 \(c\) 层里 \(0<p(x\given c)<1\)（否则那一层根本没有可比的两组人）。第三条就是后面要单独讲的正值性。

这是充分条件，不是唯一路线。同一张图上常有多个合法调整集，它们估出的是同一个量，方差和测量成本却相差很远。也有些图里任何协变量集合都不够用，得靠前门准则、工具变量或更一般的 do-演算——那类问题不是``再多控制几个变量''能解决的。

\subsection{分层内一致而汇总相反，差别在人群构成}

设 \(C=0\) 表示轻症，\(C=1\) 表示重症。某治疗在两层内都把康复率抬高 \(10\) 个百分点：

\[
\begin{array}{c|cc}
& X=1 & X=0\\
\hline
C=0 & 0.90 & 0.80\\
C=1 & 0.30 & 0.20
\end{array}
\]

若治疗组有 \(90\%\) 是重症，对照组只有 \(10\%\) 是重症，则两组的总体康复率分别是

\[
0.1\times0.90+0.9\times0.30=0.36,
\qquad
0.9\times0.80+0.1\times0.20=0.74.
\]

每层都更好，汇总起来却更差。数字没有矛盾，两个平均用的是不同的权重：治疗组按自己的病情构成加权，对照组按自己的。调整公式要求两边都按目标总体的 \(p(c)\) 加权，那样算出来的差就回到 \(10\) 个百分点。

但也不能就此宣布``分层的答案永远对''。这个结论成立，前提是 \(C\) 确实是处理前的共同原因。如果 \(C\) 是治疗引发的中间状态，分层就成了控制中介，削掉的是效应本身。该汇总还是该分层，图和研究设计说了算，数据表本身不说话。

\subsection{条件在共同结果上会凭空造出关联}

碰撞点是路径上两个箭头共同指向的那个节点。在

\[
X\longrightarrow S\longleftarrow Y
\]

里，不给定 \(S\) 时这条路径是关着的，\(X\) 与 \(Y\) 无关；一旦只分析 \(S=1\) 的样本，或者在回归里控制住 \(S\)，两者就相关起来。

设能力 \(X\) 和家庭资源 \(Y\) 在总体中独立，二者都提高录取概率 \(S\)。在已录取的学生里，能力偏弱的那些多半得靠更强的资源补上，于是样本内 \(X\) 与 \(Y\) 呈负相关。这个负号不说明能力削弱资源，它是选择规则的产物。这套机制的概率论根子在\hyperref[sec:05-Probability/02-Conditional-Probability/04]{``条件独立与信息结构''}一节，那里说清了独立与条件独立互不蕴含的两个方向。

值得警惕的是，碰撞点常常不以协变量的形式出现，而是藏在样本怎么进来的规则里。医院病例库、平台活跃用户、``完成随访者''，都是由多个因素共同决定的选择变量。碰撞点的后代同样会开路，所以缺失指示、纳入标准、存活到研究终点这些条件都该画进图里，而不是只审查数据表里现成的那几列。开头那位分析师加的``是否完成随访''，很可能正是这样一个碰撞点。

\subsection{中介变量回答的是另一个问题}

若图是

\[
X\longrightarrow M\longrightarrow Y,
\qquad
X\longrightarrow Y,
\]

则 \(M\) 是处理之后产生的中介。估总效应时控制 \(M\)，等于把经 \(M\) 传递的那部分掐掉，剩下的既不是总效应也不一定是直接效应。就算目标本来就是直接效应，控制中介也不自动正确：\(M\) 与 \(Y\) 之间可能存在被处理诱发的混杂，而自然直接效应涉及跨世界的反事实，需要比后门准则更强的假设。用药后第七天的化验指标就是典型的中介，把它放进模型，估的已经不是那个药的总效应了。

\begin{center}
\begin{tikzpicture}[>=stealth,
  var/.style={draw,circle,minimum size=0.75cm,fill=blue!6},
  every node/.style={font=\small}]
  \begin{scope}[xshift=-4cm]
    \node[var] (c) at (0,1.35) {\(C\)};
    \node[var] (x) at (-1,0) {\(X\)};
    \node[var] (y) at (1,0) {\(Y\)};
    \draw[->] (c) -- (x);
    \draw[->] (c) -- (y);
    \draw[->,thick] (x) -- (y);
    \node at (0,-0.85) {混杂：通常应控制 \(C\)};
  \end{scope}
  \begin{scope}
    \node[var] (x2) at (-1,1.35) {\(X\)};
    \node[var] (y2) at (1,1.35) {\(Y\)};
    \node[var] (s) at (0,0) {\(S\)};
    \draw[->] (x2) -- (s);
    \draw[->] (y2) -- (s);
    \node at (0,-0.85) {碰撞点：控制 \(S\) 会开路};
  \end{scope}
  \begin{scope}[xshift=4cm]
    \node[var] (x3) at (-1.2,0.65) {\(X\)};
    \node[var] (m) at (0,0.65) {\(M\)};
    \node[var] (y3) at (1.2,0.65) {\(Y\)};
    \draw[->] (x3) -- (m);
    \draw[->] (m) -- (y3);
    \draw[->,thick] (x3) to[bend left=35] (y3);
    \node at (0,-0.85) {中介：控制 \(M\) 会删掉路径};
  \end{scope}
\end{tikzpicture}
\end{center}

\emph{三个图中变量都可能与结果相关，但调整含义相反；该不该控制取决于箭头角色，而不是相关系数大小。}

三张图放在一起，能看清那条直觉错在哪。回归拿到的是变量和结果的相关强度，而这三张图里的 \(C\)、\(S\)、\(M\) 都可以与 \(Y\) 强相关；决定该不该控制的是箭头怎么走，那是数据表里没有的信息。``先把所有与 \(Y\) 相关的列放进模型''因此不是稳健策略，而是把三种角色混作一谈。正则化能在候选变量里压系数，但它压的是预测误差，认不出谁是共同原因、谁是碰撞点、谁是中介。

\subsection{调整集要在看结果之前讲清楚}

一份可信的报告应当说明：每个关键变量凭什么算处理前变量，它堵的是哪条路径，会不会反过来受处理影响，样本纳入是否依赖处理和结果。这些话必须在拿到估计值之前写下来——事后按数字挑调整集，等于用结果反选假设，置信区间随即失去含义。若几张同样合理的 DAG 导出不同的调整集，正当做法是分别算、做敏感性分析，而不是把其中一张说成事实。

因果图证明不了没有隐藏混杂。它能做的是把这个无法用当前数据检验的假定摆到明处，让人指着某个节点问``你确定这里没有第三个变量吗''。比起一句含糊的``相关变量都已控制''，这是实打实的进步。

后门准则解决的是``该控制谁''。但它依赖一张图，而图往往有争议。下一节换一套语言：不先画图，直接把处理效应定义在个体的两个世界之间，看看需要哪些假设才能把它算出来。
